% Transcription — fonds Alexandre Grothendieck, shelfmark 151, batch 1 (pages 1-20).
% Established from the facsimile archives/batches/151/batch-01.pdf (a mirror of
% https://grothendieck.umontpellier.fr/151.pdf), per the skill
% transcribe-grothendieck.
%
% Pages 9 and 15 are blank — absent. Pages 3, 5 and 7 are three copies of one
% administrative letter (Inspection académique de Vaucluse, Avignon, 6 December
% 1990), scanned upside down and unrelated to the notes — absent. Pages 8, 17,
% 18 and 19 are written on the *back* of that same letterhead, and those versos
% do carry the mathematics, so they are transcribed like any other page.
%
% The batch holds four distinct pieces of writing, in this order: a table of
% contents for a text called « Structures stratifiées » (page 1); four sheets of
% relation-calculus for a 1/2-simplicial system E_0, E_1, E_2, numbered 1 to 4
% by the author in circles (pages 2, 4, 6, 8); a continuous five-page text,
% « Hyperrecouvrements », paginated 1 to 5 by the author (pages 10-14); and a
% second numbered series, 1 to 3, on the system (E_n, d_n, k_n) and on
% X_1, X_2, X_3 (pages 16, 18, 19), with a sheet of lim-computations for phi_!
% between them (page 17) and three surviving lines on page 20.
%
% The dating is the inventory's own and spans the group's range; the notes
% themselves carry no date. The 1990 letter reused as paper is later than the
% « (1981) » of the folder title, which is consistent with the inventory
% proposing a range rather than a year — but the pass does not conclude that,
% and the observation stays out of \dating{}.
%
% Pass: Opus 5 (claude-opus-5), 2026-08-15 — first pass, unchecked against the pages by a human.
\documentclass[11pt,a4paper]{article}
% Le préambule commun à toutes les transcriptions du fonds Grothendieck.
%
% Il tient en une idée : **l'incertitude se compose, elle ne se résout pas.**
% Une transcription qui lisse un mot douteux a détruit la seule chose pour
% laquelle on la faisait — un lecteur doit pouvoir savoir, à chaque endroit,
% ce qui était sur le feuillet et ce qui a été ajouté. Les six macros qui
% suivent sont donc l'appareil critique, et non de la décoration.
%
% Les mêmes commandes sont reconnues par `scripts/render.mjs`, qui produit la
% vue de lecture du site. Une macro ajoutée ici doit l'être aussi là-bas,
% faute de quoi le rendu échouera bruyamment — ce qui est voulu : un rendu
% silencieusement faux est pire qu'un rendu absent.

\ProvidesPackage{grothendieck}[2026/08/08 Transcriptions du fonds Grothendieck]

\usepackage{amsmath,amssymb,amsthm}
\usepackage{mathrsfs}
\usepackage{stmaryrd}
% Les diagrammes commutatifs sont partout dans ce fonds : c'est la langue
% même de la géométrie algébrique de Grothendieck, et une transcription qui
% les rendrait en prose ne transcrirait rien.
\usepackage{tikz-cd}
% Français par défaut — c'est la langue des feuillets — mais l'anglais est
% chargé aussi : la traduction et le résumé sont des documents anglais, et la
% typographie française (espaces avant les deux-points) y serait une faute.
% Ces documents basculent par \selectlanguage{english} en tête de corps.
\usepackage[english,french]{babel}
\usepackage{fontspec}
\usepackage{geometry}
\geometry{a4paper,margin=25mm,marginparwidth=28mm}
\usepackage{marginnote}
\usepackage{xcolor}
% ulem, not soul. `soul`'s \st and \ul are built on a character-by-character
% scan that XeLaTeX's font machinery breaks: the document fails with an
% undefined control sequence rather than degrading. ulem does the same two jobs
% and is Unicode-engine safe. `normalem` stops it hijacking \emph, which the
% transcriptions use for Grothendieck's own emphasis.
\usepackage[normalem]{ulem}
\usepackage{hyperref}
\hypersetup{colorlinks=true,linkcolor=black,urlcolor=black,pdfborder={0 0 0}}

% --- Métadonnées du lot -----------------------------------------------------
% Reprises telles quelles par le script de rendu, qui les affiche en tête de
% la vue de lecture. Elles fixent ce que le fichier prétend couvrir : sans
% elles, une transcription flotte sans point d'attache dans un fonds de seize
% mille feuillets.

\newcommand{\folder}[1]{\gdef\grothFolder{#1}}
\newcommand{\batch}[1]{\gdef\grothBatch{#1}}
\newcommand{\leaves}[2]{\gdef\grothFirst{#1}\gdef\grothLast{#2}}
\newcommand{\foldertitle}[1]{\gdef\grothFoldertitle{#1}}
\gdef\grothFolder{?}\gdef\grothBatch{?}\gdef\grothFirst{?}\gdef\grothLast{?}\gdef\grothFoldertitle{}

% --- Le résumé --------------------------------------------------------------
%
% Il ouvre la lecture modernisée, sous le titre « Résumé », et il est composé
% en petit. Ce n'est pas un ornement : c'est le seul passage du document qui
% ne soit pas des mathématiques, et un lecteur doit pouvoir le reconnaître
% avant de l'avoir lu — puis le sauter s'il connaît déjà le sujet, ou s'y
% arrêter s'il ne le connaît pas. Le filet qui le suit marque la reprise du
% corps.

\newenvironment{resume}
  {\small\setlength{\parskip}{0.45em}}
  {\par\medskip\hrule\medskip}

% --- Mots-clés ---------------------------------------------------------------
%
% En anglais, en clôture du résumé de la lecture modernisée : le vocabulaire
% moderne sous lequel chercher ce que les feuillets construisent. Le manifeste
% du site (`npm run manifest`) les extrait pour étiqueter la cote entière —
% cette ligne est la source unique des tags, il n'y a pas de fichier à côté.
\newcommand{\keywords}[1]{\par\smallskip\noindent
  {\footnotesize\textsc{Keywords} --- \textit{#1}}\par}

% --- L'appareil critique ----------------------------------------------------

% Le numéro de feuillet, en marge. C'est l'ancre qui permet de régler un
% désaccord sur une formule en regardant *un* feuillet plutôt qu'une chemise
% de six cents. Le site s'en sert aussi pour tourner les pages du fac-similé
% quand on fait défiler la transcription.
\newcommand{\leaf}[1]{%
  \marginnote{\footnotesize\color{gray!70}\textsf{#1}}%
  \label{leaf:#1}\ignorespaces}

% Illisible : on ne devine pas. Un mot inventé qui se lit comme les autres est
% le pire résultat possible.
\newcommand{\ill}{\textcolor{red!60!black}{[\,\ldots\,]}}

% Lecture douteuse : le mot est donné, et signalé comme incertain.
\newcommand{\uncertain}[1]{\uline{#1}}

% Ajout de l'éditeur — une lettre manquante, un mot sous-entendu.
\newcommand{\add}[1]{\textcolor{blue!50!black}{[#1]}}

% Biffé par Grothendieck lui-même. Ce qu'il a rayé fait partie du document :
% une rature dit souvent où la pensée a changé de direction.
\newcommand{\struck}[1]{\sout{#1}}

% Note du transcripteur, sur l'état matériel du feuillet.
\newcommand{\note}[1]{\par\smallskip{\small\color{gray!60!black}\textit{[#1]}}\par\smallskip}

% Note marginale de Grothendieck, distincte du corps du texte.
\newcommand{\marginal}[1]{\par\smallskip\noindent\hspace{1em}%
  {\small\color{gray!60!black}#1}\par\smallskip}

% --- Titre ------------------------------------------------------------------

\AtBeginDocument{%
  \begin{center}
    {\large\bfseries Fonds Alexandre Grothendieck --- cote n\textsuperscript{o}~\grothFolder}\\[2pt]
    {\itshape\grothFoldertitle}\\[4pt]
    {\small feuillets \grothFirst--\grothLast\ (lot \grothBatch)}\\[2pt]
    {\footnotesize\color{gray!60!black}Transcription établie d'après le fac-similé numérisé
     par l'Université de Montpellier.\\
     Les lectures incertaines sont soulignées, les illisibles notées [\,\ldots\,],
     les ajouts entre crochets.}
  \end{center}
  \vspace{1em}}

\endinput


\folder{151}
\batch{1}
\pages{1}{20}
\dating{[à partir de 1981-à partir de 1990]}
\watermark{Édition de démonstration}
\foldertitle{Topos stratifié (1981) : notes manuscrites (s.d.).}

\begin{document}

\section*{« Structures stratifiées » — le plan}

\page{1}
Structures stratifiées

\begin{enumerate}
  \item La situation type — préliminaires d'un programme. \quad 1.
  \item Stratifications globales : le formalisme $\tfrac12$ simplicial (sans tubes). \quad 10
  \item \struck{\ill} Introduction des tubes. \quad 15
  \item Topos canoniques associés à une stratification (globale). \quad 23
\end{enumerate}

\note{Les nombres qui ferment chaque ligne sont portés dans la marge droite du
sommaire : ce sont les renvois de pagination du texte annoncé, et non les
numéros de l'archiviste. Aucune des pages du présent lot ne porte cette
pagination-là, sauf le texte « Hyperrecouvrements » ci-dessous, paginé 1 à 5
par l'auteur.}

\note{Au point 3, un mot est biffé si lourdement qu'il ne se lit plus ; les
premières lettres suggèrent une majuscule, nous n'allons pas plus loin.}

\section*{Le formalisme $\tfrac12$-simplicial : le système $E_0, E_1, E_2$}

\note{Les quatre pages qui suivent sont numérotées 1 à 4 par l'auteur, dans un
cercle en haut de page ; ce sont des feuilles de travail, où les relations sont
écrites, encadrées, biffées et reprises. Nous transcrivons les systèmes de
relations et les phrases, sans reproduire chaque reprise : là où un fragment est
recouvert au point de ne plus se lire, il est noté. Les accolades que l'auteur
place sous un terme pour le renommer sont conservées telles quelles : ainsi
$\underbrace{\rho\delta_1}_{\delta_3}$.}

\page{2}
Le système est celui de trois ensembles et des applications qui les relient :

\begin{tikzcd}[column sep=large]
E_0 \arrow[r, "\delta_0", bend left=18] & E_1 \arrow[l, "p_0"', bend left=18] \arrow[r, "\delta_1", bend left=18] & E_2 \arrow[l, "p_{12}"', bend left=18]
\end{tikzcd}

outre $p_0$ et $p_{12}$, le système comporte
$p_1 \colon E_1 \to E_0$, puis $p_{23}, p_{31} \colon E_2 \to E_1$, et les
faces $\delta_2, \delta_3 \colon E_1 \to E_2$ : toutes sont nommées dans les
relations ci-dessous.

$\mathfrak{S}_2$ opère sur $E_1$, d'élément $\sigma$ ; $\mathfrak{S}_3$ opère sur
$E_2$, d'éléments $\rho$, $\sigma_1$, $\sigma_2$, $\sigma_3$, $\rho^2$.

\note{Le coin supérieur gauche de la page porte, sous ce diagramme, un faisceau
de flèches recourbées reliant $E_1$ à $E_2$ et étiquetées $\delta_1$,
$\delta_2$, $\delta_3$, en partie recouvert de hachures. Nous ne dessinons
qu'une flèche par sens, faute de pouvoir en séparer trois lisiblement ; aucune
n'est perdue, les autres sont nommées à la ligne suivante et opèrent dans les
relations.}

Les opérations sont données en coordonnées :

\[
  \delta_3(x,y) = (x,x,y), \qquad
  \delta_1(x,y) = (y,x,x), \qquad
  \delta_2(x,y) = (x,y,x),
\]
\[
  p_{12}(x,y,z) = (x,y), \qquad
  p_{23}(x,y,z) = (y,z), \qquad
  p_{31}(x,y,z) = (z,x),
\]
\[
  \rho(x,y,z) = (y,z,x).
\]

En marge de la première boîte : $\delta_1\sigma(x,y) = \uncertain{(y,y,x)}$,
puis $\delta_3\delta_3$.

$p_1 = p_0\sigma$, $p_0 = p_1\sigma$.

\struck{$p_{12} = p_{23}$}

\[
  \begin{cases}
    p_{23} = p_{12}\rho \\
    p_{31} = p_{23}\rho = p_{12}\rho^2 \\
    p_{12} = p_{31}\rho = p_{23}\rho^2
  \end{cases}
  \qquad
  \begin{cases}
    \delta_1 = \rho\delta_2 = \rho^2\delta_3 \\
    \delta_2 = \rho\delta_3 = \rho^2\delta_1 \\
    \delta_3 = \rho\delta_1 = \rho^2\delta_2
  \end{cases}
\]
\[
  \sigma p_{12} = p_{12}\sigma_3, \qquad
  \sigma p_{23} = p_{23}\sigma_1, \qquad
  \sigma p_{31} = p_{31}\sigma_2.
\]

\uncertain{Il} suffit de se donner $p_0, p_{12}, \delta_0, \delta_1$ avec les
trois ensembles $E_0$, $(E_1,\mathfrak{S}_2)$, $(E_2,\mathfrak{S}_3)$, avec les
relations

\[
  \boxed{p_0\delta_0 = \mathrm{id}_{E_0}}, \qquad \boxed{\sigma\delta_0 = \delta_0},
\]
\[
  \delta_0p_0 = p_{12}\,\underbrace{\rho\delta_1}_{\delta_3}
             = p_{23}\,\delta_1
             = p_{31}\,\underbrace{\rho^2\delta_1}_{\delta_2}.
\]

\note{Cette dernière chaîne d'égalités est écrite deux fois sur la page, la
seconde fois précédée d'un « i.e. », et les termes intermédiaires y sont
surchargés d'accolades qui renomment $\rho\delta_1$ en $\delta_3$ et
$\rho^2\delta_1$ en $\delta_2$. Une inégalité barrée la traverse, flanquée de
deux points d'interrogation :}

\[
  (p_{12}\rho)\delta_1 \neq p_{23}(\rho\delta_1) \qquad ??
\]

\note{Le désaccord porte sur l'ordre d'application, et il est tranché plus bas :
la boîte suivante donne les trois valeurs, qui ne sont pas les mêmes selon le
$\delta$ auquel on compose.}

\[
  \begin{cases}
    \delta_0p_0 = p_{12}\,\underbrace{\rho\delta_1}_{\delta_3} \\
    \mathrm{id}_{E_1} = p_{12}\,\underbrace{\rho^2\delta_1}_{\delta_2} \\
    \sigma = p_{12}\,\delta_1
  \end{cases}
\]
\[
  \boxed{\sigma_3\delta_1 = \delta_1}, \qquad \boxed{p_{12}\sigma_3 = \sigma p_{12}},
\]
\[
  \boxed{\rho\,\delta_1\delta_0 = \sigma_1\,\delta_1\delta_0 = \delta_1\delta_0}, \qquad
  p_0\,p_{12} = p_0\,\underbrace{p_{12}\rho^2}_{p_{31}}.
\]

\page{4}
En tête de page, les générateurs sont recensés en trois colonnes —
$p_0/p_1$ et $\delta_0$ ; $p_{12}, p_{13}, p_{31}$ au-dessus de
$p_{21}, p_{32}, p_{23}$ ; $\delta_1, \delta_2, \delta_3$ — puis
$\delta_0p_0, \delta_0p_1, \mathrm{id}, \sigma$. Trois longs traits obliques
traversent le recensement.

\begin{tikzcd}[column sep=large]
E_0 \arrow[r, "\delta_0", bend left=18] & E_1 \arrow[l, "p_0", bend left=18] \arrow[r, "\delta_1"', bend right=18] & E_2 \arrow[l, "p_{12}"', bend right=18]
\end{tikzcd}

\note{L'auteur inscrit $\mathfrak{S}_2$ au-dessus de $E_1$ et $\mathfrak{S}_3$
au-dessus de $E_2$ : c'est la donnée minimale annoncée à la page précédente,
redessinée seule dans un cadre.}

\[
  \boxed{p_0\delta_0 = \mathrm{id}_{E_0}}, \qquad \boxed{\sigma\delta_0 = \delta_0},
\]
\[
  \begin{cases}
    p_{12}\delta_1 = \sigma & \text{(involution)} \\
    p_{12}\,\underbrace{\rho\delta_1}_{\delta_3} = \delta_0p_0 & \text{(idempotent)} \\
    p_{12}\,\underbrace{\rho^2\delta_1}_{\delta_2} = \mathrm{id} &
  \end{cases}
\]
\[
  \boxed{\sigma_1\delta_1 = \delta_1}, \qquad
  \boxed{p_{12}\sigma_3 = \underbrace{\sigma\,p_{12}}_{p_{21}}},
\]
\[
  \boxed{\rho(\delta_1\delta_0) = \sigma_1(\delta_1\delta_0) = \delta_1\delta_0}
  \quad \text{i.e.} \quad g(\delta_1\delta_0) = \delta_1\delta_0 \ \text{ pour } g \in \mathfrak{S}_3,
\]
\[
  \boxed{p_0\,p_{12}\,\sigma_3 = \underbrace{p_0\,p_{12}}_{p_1}}.
\]

\note{Quatre listes entourées occupent la marge droite : $p_0, p_1$ ;
$\delta_1, \delta_2, \delta_3$ ; $p_{12}, p_{23}, p_{31}, p_{21}, p_{32},
p_{13}$ ; $g_1, g_2, g_3$ — et, contre la première, les définitions
$\delta'_1 = \delta_1\sigma$, $\delta'_2 = \delta_2\sigma$,
$\delta'_3 = \delta_3\sigma$. Ce sont des inventaires de générateurs, non des
relations.}

\page{6}
\[
  p_0p_{12} = p_0p_{13} = p_0\,\underbrace{p_{12}\rho^2}_{p_{31}}\,\sigma_2,
  \qquad \rho^2\sigma_2 = \sigma_3\rho^2,
\]
\[
  p_0\,p_{12}\,\sigma_1 = p_0\,p_{12}.
\]

Le calcul de $\rho^2\sigma_2$ est mené en coordonnées :
$(x,y,z) \mapsto (z,y,x) \mapsto (y,x,z)$, et
$(x,y,z) \mapsto (z,x,y)$, enfin $(x,z,y) = \sigma_1$.

\note{Cette colonne de calculs est traversée de deux longues courbes qui
enferment le résultat ; les images intermédiaires sont écrites plusieurs fois,
et nous ne retenons que la chaîne complète. À gauche, $\delta_1\sigma =
\uncertain{(x,y,y)}$, avec les triplets $(x,x,y)$ et $(y,y,x)$ notés dessous.}

Si $E_0$ est réduit à un point, les $\delta_0, \delta_1, \partial_0$, $E_1$,
$E_2$ deviennent des ens. pointés, et $\mathfrak{S}_2$, $\mathfrak{S}_3$
respectent le pointage (« pointés »). \struck{\ill\ respectant le} Dans, en
plus des opérations de $\mathfrak{S}_2, \mathfrak{S}_3$, on a

\begin{tikzcd}[column sep=large]
E_1 \arrow[r, "\delta_1", bend left=18] & E_2 \arrow[l, "p_{12}"', bend left=18]
\end{tikzcd}

i.e. $\delta_1$ \uncertain{et $p_{12}$} \uncertain{respectant} les pointations,
et on a $p_{12}, \delta_1$ satisfaisant aux cinq relations supplémentaires

\[
  \begin{cases}
    p_{12}\delta_1 = \sigma \\
    p_{12}\,\underbrace{\rho\delta_1}_{\delta_3} = \text{application constante de valeur } e_1 \\
    p_{12}\,\underbrace{\rho^2\delta_1}_{\delta_2} = \mathrm{id}_{E_1} \\
    \sigma_1\delta_1 = \delta_1 \\
    p_{12}\sigma_3 = \sigma\,p_{12}
  \end{cases}
\]

\note{Une seconde boîte, à droite de celle-ci, refait le même compte à partir de
$\delta_2$ au lieu de $\delta_1$ : $p_{12}\delta_2 = \mathrm{id}_{E_1}$,
$p_{12}\rho\delta_2 = \sigma$, $p_{12}\rho^2\delta_2 = e_1$ — l'application
constante — avec $\sigma_2\delta_2 = \delta_2$ et $p_{12}\sigma_3 = \sigma
p_{12}$. Elle est en partie biffée, et une partie des indices y est
illisible ; le lecteur y trouvera la même liste de cinq relations, transportée
par $\rho$.}

\page{8}
\[
  p_{12}\delta_1(x) = \sigma(x), \qquad
  \underbrace{p_{12}\,\rho\delta_1}_{p_{23}}(x) = \delta_0p_0(x), \qquad
  \underbrace{p_{12}\,\rho^2\delta_1}_{p_{31}}(x) = x.
\]

\marginal{NB $\delta_1(x)$ ne peut être fixé par $\rho$ (donc par
$\mathfrak{S}_3$) que si $x = \delta_0p_0(x)$ (condition néc. et suff.).}

\[
  \begin{cases}
    \delta_0p_0(x) = 1 \\
    \sigma(x)\cdot x = 1 \\
    p_{12}(z)\,p_{23}(z)\,p_{31}(z) = 1
  \end{cases}
  \qquad
  \sigma(x)\,\underbrace{\delta_0p_0(x)}_{=\,1}\,x = 1
\]

pour $z \in E_2$ (une relation \struck{\ill} supplémentaire par orbite sous
$\mathfrak{S}_3$ \struck{qui ne rencontre pas} $\delta_1(E_1)$).

\note{Les trois relations encadrées sont écrites multiplicativement, comme des
relations de groupe : c'est le seul endroit du lot où le système
$\tfrac12$-simplicial est lu comme une présentation.}

Pour $\alpha \in E_1$, avec $F_\alpha \subset F \times F$ au-dessus de $F$ :

\begin{tikzcd}[column sep=small, row sep=small, nodes={font=\scriptsize}]
F \times F \times F \arrow[d, "p_{12}", bend left=18] \arrow[d, "p_{13}"', bend right=18] \\
F \times F
\end{tikzcd}

\[
  E_2(\alpha) \ni \delta_1(\alpha) = \delta_1\sigma(\alpha),
\]
\[
  \Pi_0(F_\alpha \times_F F_\alpha) = \underbrace{p_{12}^{-1}(\alpha) \cap p_{13}^{-1}(\alpha)}_{E_2(\alpha)},
  \quad \text{stable par } \sigma_1.
\]

Prenons $\alpha \neq \delta_1$, alors $\delta_2 \notin p_{12}^{-1}(\alpha)$ ;
les pts fixes par $\sigma_1$ sont dans image $\delta_1(E_1)$, mais
$p_{12}(\delta_1(x)) = p_{13}(\delta_1(x)) = \sigma x$, donc
\struck{\ill} $\delta_1(x) \in E_2(\alpha) \iff \sigma x = \alpha$, i.e.
$x = \sigma\alpha$, et $\delta_1\sigma(\alpha)$ \uncertain{est le seul} des
$\delta_2 = \delta_1(\alpha)$.

\note{La dernière proposition est écrite par-dessus une rature et le verbe ne se
lit qu'à demi ; le sens — que $\delta_1\sigma(\alpha)$ est le seul point fixe
en jeu — est celui que porte la ligne, mais la lecture du mot reste douteuse.}

\section*{Hyperrecouvrements}

\note{Ce qui suit est un texte suivi, d'une seule venue, paginé 1 à 5 par
l'auteur en haut de page ; nous le donnons intégralement. C'est la seule partie
du lot qui soit rédigée plutôt que calculée.}

\page{10}
\note{Page 1 de l'auteur.}

\textbf{Hyperrecouvrements.}

Soit $X$ une catégorie $U$-topos. On veut définir \struck{une} la notion
d'\emph{hyperrecouvrement} de $X$ (ou encore, du topos défini par $X$), indexé
par une petite catégorie $A$. \struck{Un tel type} \uncertain{C'est} par
définition un foncteur

\[
  \varphi \colon A \longrightarrow X, \qquad
  \alpha \longmapsto U_\alpha\ (= \varphi(\alpha)),
\]

satisfaisant certaines conditions, inspirées par les \ill\ des
« \uncertain{Čechistes} » recouvrements dont les nerfs, et le raffinement de ce
\ill\ suggéré par Cartier (pour les espaces topologiques ou paracompacts).
Voici les conditions qui me viennent à l'esprit :

\note{La même abréviation de deux lettres, ponctuée, revient aux deux endroits
marqués \ill\ — « les \ill\ des recouvrements », puis « ce \ill\ suggéré par
Cartier ». Elle désigne un procédé, et non un objet, mais nous ne la lisons pas
et ne la devinons pas. Le mot entre guillemets est inséré au-dessus de la ligne
par un signe d'insertion.}

\medskip
\textbf{Hyp 1} \quad Les objets $U_\alpha$ ($\alpha \in \mathrm{Ob}\,A$)
recouvrent l'objet final $e_X$ de $X$.

\medskip
\textbf{Hyp 1${}'$} \quad Pour tout $\alpha_0 \in \mathrm{Ob}\,A$, le
foncteur induit $A/\alpha_0 \to X/U_{\alpha_0}$ satisfait à Hyp 1, i.e. l'ens.
des flèches $U_\alpha \to U_{\alpha_0}$, associées aux flèches
$\alpha \to \alpha_0$ de but $\alpha_0$ dans $A$, \page{11} soit couvrant.

\note{Hyp 1${}'$ est barrée d'une grande croix qui court sur les quatre
dernières lignes de la page 10, et une note de quelques mots, écrite en travers
dans le coin inférieur gauche, commence par « Est-ce que » ; elle est trop
recouverte pour être lue. L'énoncé lui-même reste lisible et sert à la suite —
il est repris tel quel à la page 3 de l'auteur.}

\medskip
\textbf{Hyp 2} \quad Pour $\alpha, \beta \in \mathrm{Ob}\,A$, l'ens. des
\struck{flèches} diagrammes $\gamma \to \alpha$, $\gamma \to \beta$ dans $A$
(i.e. l'ens. des objets de $A/\alpha\times\beta$, où $\alpha\times\beta$ est le
produit cartésien dans $\widehat{A}$), indexant les flèches correspondantes
$U_\gamma \to U_\alpha \times U_\beta$, donne une famille couvrante
$U_\alpha \times U_\beta$.

\medskip
\marginal{[Hyp 2 généralisé :] Idem pour une suite finie
$(\alpha_0, \alpha_1, \dots, \alpha_n)$ ($n \in \mathbb{N}^*$) d'objets de $A$,
et les \struck{diagr.} flèches
$U_\gamma \to U_{\alpha_0} \times U_{\alpha_1} \times \dots \times U_{\alpha_n}$
indexées par les objets de
$A/\alpha_0\times\alpha_1\times\dots\times\alpha_n$.}

Pour $n = 0$, on retrouve Hyp 1${}'$, pour $n = 1$, la formulation initiale de
Hyp 2. Mais Hyp 2 pour $n = 1$ implique-t-il Hyp 2 pour \struck{\ill}
$n \geq 1$ ? On voit qu'il en est ainsi, par récurrence, en utilisant le fait
que la famille composée de deux familles couvrantes

\begin{tikzcd}[column sep=large, nodes={font=\scriptsize}]
U_{\gamma'} \arrow[r, "\varphi(v)"] & U_\gamma \times U_{\alpha_{n+1}} \arrow[r, "\varphi(u)\times\mathrm{id}"] & (U_{\alpha_0} \times \dots \times U_{\alpha_n}) \times U_{\alpha_{n+1}}
\end{tikzcd}

($u$ parcourant les flèches $\gamma' \to \gamma \times \alpha_{n+1}$, et $v$ les
flèches $\gamma \to \alpha_0 \times \dots \times \alpha_n$) est couvrante. Ainsi
Hyp 1${}'$ \struck{\ill} Hyp 2.

\note{Les deux étiquettes de ce diagramme sont écrites très petit au-dessus des
flèches ; nous lisons $\varphi(v)$ sur la première et $\varphi(u)\times
\mathrm{id}$ sur la seconde, mais la phrase qui suit assigne à $u$ le rôle de la
première et à $v$ celui de la seconde. L'un des deux est mal placé, sur la page
ou sous notre lecture ; nous ne corrigeons pas. Le mot biffé de la dernière
ligne portait vraisemblablement une implication, remplacée par un signe
surchargé.}

\page{12}
\note{Page 3 de l'auteur.}

… forme initiale (avec $\alpha, \beta$) implique Hyp 2 sous sa forme
générale (avec termes en $\alpha_0, \alpha_1, \dots \alpha_n$).

\medskip
\textbf{Hyp 2${}'$} : Pour tout $\alpha_0 \in \mathrm{Ob}\,A$, le
foncteur induit $A/\alpha_0 \twoheadrightarrow X/U_{\alpha_0}$ satisfait :
Hyp 2, i.e. pour tout diagramme $\alpha_1 \to \alpha_0 \leftarrow \alpha_2$,
regardant les couples commutatifs de $A$ qui les complètent

\begin{tikzcd}[column sep=small, row sep=small]
& \gamma \arrow[dl] \arrow[dr] & \\
\alpha_1 \arrow[dr] & & \alpha_2 \arrow[dl] \\
& \alpha_0 &
\end{tikzcd}

i.e. les objets de $A/\alpha_1\times_{\alpha_0}\alpha_2$, la famille des flèches
correspondantes

\[
  U_\gamma \longrightarrow U_{\alpha_1} \times_{U_{\alpha_0}} U_{\alpha_2}
\]

est couvrante (i.e., épimorphique).

Les hypothèses Hyp 1, 2, 2${}'$ se reformulent plaisamment, en termes du
prolongement canonique

\[
  \varphi_! \colon \widehat{A} \longrightarrow X
\]

de $\varphi$ qui en un foncteur \marginal{$\widehat{A} \to X$} commutant et aux
petites $\varinjlim$. Ce sont en particulier \page{13} \note{Page 4 de
l'auteur.} les conditions :

\begin{enumerate}
  \item[a)] $\varphi_!(e_{\widehat{A}}) \longrightarrow e_X$ est épi
  \item[b)] $\varphi_!(\alpha\times\beta) \longrightarrow \varphi_!(\alpha)\times\varphi_!(\beta)$ est épi, $\forall\,\alpha,\beta$ dans $\mathrm{Ob}\,A$
  \item[c)] $\varphi_!(\alpha\times_{\alpha_0}\beta) \longrightarrow \varphi_!(\alpha)\times_{\varphi_!(\alpha_0)}\varphi_!(\beta)$ est épi, $\forall$ diagramme $\alpha \to \alpha_0 \leftarrow \beta$ dans $A$.
\end{enumerate}

Un cas particulièrement agréable où ces conditions sont satisfaites, c'est quand
$\varphi_!$ est exact : \uncertain{gauche}, i.e. qu'il correspond \struck{\ill}
à un morphisme de topos

\[
  f \colon X \longrightarrow \widehat{A}
\]

par la formule $\varphi_! = f^*$.

Alors dans a) b) c) les flèches épi en question sont même des iso, et
\uncertain{je puis} que l'inverse est vrai également. Quand — il est
essentiellement \uncertain{suivi}, en écrivant \uncertain{conséquemment}, pour un
$F \in \mathrm{Ob}\,\widehat{A}$, $F = \varinjlim_{A/F}\alpha$ \ill.
On voit que \struck{\ill} b) et c)

\note{Les deux dernières lignes de la page sont couvertes par une longue rature
horizontale qui emporte le verbe principal ; ce qui reste lisible est donné
ci-dessus, et les mots que nous ne tenons pas pour sûrs sont soulignés. La suite
de la phrase se reprend au haut de la page suivante.}

\page{14}
\note{Page 5 de l'auteur.}

… impliquent

\begin{enumerate}
  \item[b${}'$)] $\varphi_!(F\times G) \longrightarrow \varphi_!(F)\times\varphi_!(G)$ épi, $\forall\,F,G$ dans $\mathrm{Ob}\,\widehat{A}$
  \item[c${}'$)] $\varphi_!(G\times_F H) \longrightarrow \varphi_!(G)\times_{\varphi_!(F)}\varphi_!(H)$ épi, $\forall$ diagramme $G \to F \leftarrow H$ dans $\widehat{A}$.
\end{enumerate}

et de là, plus généralement, la forme équivalente unique de Hyp 1, 2, 2${}'$ :

\medskip
\textbf{(Hyp)} \quad Pour tout diagramme fini $(F_i)_{i \in I}$ dans
$\widehat{A}$, la flèche canonique

\[
  \varphi_!\Bigl(\varprojlim_I F_i\Bigr) \longrightarrow \varprojlim_I \varphi_!(F_i)
\]

est épi.

(Il suffit d'ailleurs de le postuler quand les $F_i$ sont dans $A$, et dans
trois cas particuliers $I = {\cdot}$, ${\cdot}\ {\cdot}$,
$({\cdot} \to {\cdot} \leftarrow {\cdot})$, correspondants aux axiomes Hyp 1,
Hyp 2, Hyp 2${}'$.)

Mais cette hypothèse peut-elle être vérifiée, sans que $\varphi_!$ soit
$=$ exact ? Quand ?

\section*{Le système $(E_n, \partial_n, k_n)$ et son dual}

\note{Seconde série numérotée par l'auteur, 1 à 3, dans un cercle en haut de
page : pages 16, 18 et 19. La page 17, intercalée, poursuit le calcul de
$\varphi_!$ du texte précédent ; nous la laissons à sa place, qui est celle du
papier.}

\page{16}
\begin{tikzcd}[column sep=large]
E_n \arrow[r, "\partial_n", bend left=18] & E_{n+1} \arrow[l, "k_n", bend left=18]
\end{tikzcd}

$n \in \mathbb{N}^*$ i.e. $n \geq 1$.

\[
  \partial_n \colon [1,n] \hookrightarrow [1,n+1] \quad \text{inclusion}, \qquad
  k_n \colon [1,n] \longleftarrow [1,n+1] \quad \text{rétraction}.
\]

$\mathfrak{S}_n$ opère sur $E_n$ ; $\mathfrak{S}_n \overset{i_n}{\hookrightarrow} \mathfrak{S}_{n+1}$.

\[
  \begin{cases}
    k_n\partial_n = \mathrm{id}_{E_n} \\
    \partial_n\sigma = i_n(\sigma)\,\partial_n & \text{si } \sigma \in \mathfrak{S}_n,\ \text{pour } n \geq 1 \\
    \sigma k_n = k_n\,i'_n(\sigma) & \text{si } \sigma \in \mathfrak{S}_{n-1},\ \text{pour } n \geq 2 \\
    k_n\sigma_{n+1} = k_n &
  \end{cases}
\]

avec $\partial_nk_n \overset{\text{déf}}{=} \rho_{n+1}$ (projecteur dans
$E_{n+1}$, rétraction sur $E_n \subset E_{n+1}$), où
$i'_n = i_n i_{n-1} \colon \mathfrak{S}_{n-1} \to \mathfrak{S}_n \to
\mathfrak{S}_{n+1}$ et où $\sigma_m \in \mathfrak{S}_m$ est la transposition
$(m-1,m)$.

Le système dual :

\begin{tikzcd}[column sep=large]
E^*_n \arrow[r, "d_n = k^*_n"', bend right=18] & E^*_{n+1} \arrow[l, "p_n = \partial^*_n"', bend right=18]
\end{tikzcd}

\[
  (*) \quad
  \begin{cases}
    \partial^*_n k^*_n = \mathrm{id}_{E^*_n} \\
    \sigma\,\partial^*_n = \partial^*_n\,i_n(\sigma) & (\sigma \in \mathfrak{S}_n) \\
    k^*_n\,\sigma = i'_n(\sigma)\,k^*_n & \\
    \sigma_{n+1}k^*_n = k^*_n &
  \end{cases}
\]

réalisation dans (Ens) de syst. des $(E^*_n, \partial^*_n, k^*_n)$
satisfaisant \ill\ $\mathfrak{S}_n$-ens. \uncertain{aux} relations $(*)$, plus

\[
  (**) \quad
  E^*_{p+q} \xrightarrow{\ (\alpha^*_{p+q,p},\ \beta^*_{p+q,q})\ } E^*_p \times E^*_q
  \quad \text{surjectif pour } p, q \geq 1.
\]

\[
  \alpha^*_{n,p} \colon E^*_n \to E^*_p \quad (n > p \geq 1), \qquad
  \boxed{\alpha^*_{n,p} = \partial^*_p\,\partial^*_{p+1} \cdots \partial^*_{n-1}}
\]
\[
  E^*_n \xrightarrow{\ \partial^*_{n-1}\ } E^*_{n-1} \xrightarrow{\ \partial^*_{n-2}\ } E^*_{n-2} \cdots \xrightarrow{\ \partial^*_p\ } E^*_p
\]
\[
  \beta^*_{n,q} \colon E^*_n \to E^*_q \quad (n \geq q \geq 1), \qquad
  \boxed{\beta^*_{n,q} = \partial'^*_q\,\partial'^*_{q+1} \cdots \partial'^*_{n-1}}
\]
\[
  E^*_n \xrightarrow{\ \partial'^*_{n-1}\ } E^*_{n-1} \xrightarrow{\ \partial'^*_{n-2}\ } E^*_{n-2} \cdots \xrightarrow{\ \partial'^*_q\ } E^*_q
\]
\[
  \partial'^*_m \colon E^*_{m+1} \to E^*_m, \qquad \boxed{\partial'^*_m = \partial^*_m\,\rho^*_{m+1}}
\]

avec $\rho^*_{m+1} = \rho_{m+1}^{-1}$, la permutation circulaire
\uncertain{directe} dans $\mathfrak{S}_{m+1}$.

Au bas de page, les versions non dualisées :
$\alpha_{n,p} \colon E_p \hookrightarrow E_n$ inclusion canonique ;
$\beta_{n,q} \colon E_q \hookrightarrow E_n$ avec $\beta_{n,q}(i) = i + (n-q)$ ;
$\partial'_n$ avec $\partial'_n(i) = i+1$ ; et
$\partial'_m = \rho_m\partial_m$ où $\rho_m \in \mathfrak{S}_{m+1}$ est donnée
par $\rho_m(i) = i+1$ si $i \in [1,m]$, $\rho_m(i) = 1$ si $i = m+1$ — la
permutation circulaire.

\page{17}
\note{Cette page, écrite au dos de la lettre administrative, reprend le calcul
de $\varphi_!$ du texte « Hyperrecouvrements ». Deux longs traits obliques la
traversent en X, sans que rien y soit remplacé : nous transcrivons ce qu'ils
recouvrent.}

\[
  F = \varinjlim_{A/F}\alpha, \qquad G = \varinjlim_{A/G}\beta,
\]
\[
  F \times G = \varinjlim_{A/F}\bigl(\alpha \times G\bigr)
             = \varinjlim_{A/F \times A/G} \alpha \times \beta,
\]
\[
  \varphi_!(F \times G) = \varinjlim_{A/F \times A/G}
    \underbrace{\varphi_!(\alpha\times\beta)}_{\varphi_!(\alpha)\times\varphi_!(\beta)},
\]
\[
  \varphi_!(F) = \varinjlim_{A/F}\varphi_!(\alpha), \qquad
  \varphi_!(G) = \varinjlim_{A/G}\varphi_!(\beta).
\]

Pour $A \xrightarrow{\ \varphi\ } X$ :

\begin{tikzcd}[column sep=large]
C = X \arrow[r, "u = \varphi^*"', bend right=18] & C' = \widehat{A} \arrow[l, "v = \varphi_!"', bend right=18]
\end{tikzcd}

\[
  u(F) = \bigl(\alpha \longmapsto \mathrm{Hom}(\varphi(\alpha), F)\bigr), \qquad
  v(G) = \varinjlim_{\alpha \in A/G} \varphi(\alpha),
\]

$u$ continu, $v$ cocontinu.

\[
  X \longrightarrow \widehat{A}, \qquad Y \mapsto \bigl(a_i \to \varphi^*(Y)\bigr), \qquad \varphi(a_i) \to Y.
\]

\note{Le haut de la page porte deux segments étiquetés $E_p$ et $E'_q$, un point
marqué $\alpha$ sur le premier, et la mention $E \setminus \{\alpha\}$ : un
dessin de la surjection $(**)$ de la page précédente, laissé sans phrase. Le
nombre porté dans la marge gauche n'est pas une pagination de l'auteur et nous
ne le reproduisons pas.}

\page{18}
cond $I, J, K \leq 3$.

\begin{tikzcd}[column sep=small, row sep=small]
I \arrow[r, "\text{épi}"] \arrow[dr, "\text{épi}"'] & J \\
& K
\end{tikzcd}

\note{Les deux flèches ci-dessus portent une double pointe : ce sont des
épimorphismes, et l'annotation « épi » est de l'auteur.}

\begin{tikzcd}[column sep=small, row sep=small]
I \arrow[r] \arrow[d] & J \arrow[r] \arrow[d] & J' \arrow[d] \\
K \arrow[r] & L \arrow[r] & M
\end{tikzcd}

\note{Les deux sommets du bas à droite ne portent pas de lettre sur la page —
le premier est biffé, le second réduit à un point ; nous les nommons $L$ et $M$
pour que le carré se dessine, et cette nomination est nôtre.}
\qquad
\begin{tikzcd}[column sep=small, row sep=small]
X_0 & X_1 \arrow[l] & X_2 \arrow[l] \\
X'_0 \arrow[u] & X'_1 \arrow[l] \arrow[u] & X'_2 \arrow[l] \arrow[u]
\end{tikzcd}

\begin{tikzcd}[column sep=small, row sep=small]
I \arrow[r, hook] \arrow[d] & J = I \amalg K \arrow[d, no head] \\
I' \arrow[r, hook] & J' = I' \amalg K
\end{tikzcd}
\qquad
\begin{tikzcd}[column sep=small, row sep=small]
X(J) & X(I \amalg K) \arrow[l] \\
X(J') \arrow[u] & X(I' \amalg K') \arrow[l]
\end{tikzcd}

\begin{tikzcd}[column sep=small, row sep=small]
E_n \arrow[r, "\partial_n"] & E_{n+1} \arrow[d, "k_{n-1}"] \\
& E_n
\end{tikzcd}

\note{La moitié basse de la page est une énumération de cas, en deux colonnes
numérotées I à V, chacune accompagnée d'un carré $E_n \to E_{n+1}$,
$E_{n-1} \to E_{n-2}$ et d'une condition sur $n$ ($n \geq 1$, $n \geq 2$,
$n \geq 3$, $n \geq 4$), avec en regard un décompte : $2, 1, 1, 1, 0$ dans la
première table ; $2, 1, 1, 1$ dans la seconde, où les lignes portent
$p = 1, q = 3$ ; $p = 2, q = 2$ ; $n = 2$ ; $n = 3$ ; $n = 4$. Le coin inférieur
droit est entièrement recouvert de ratures en boucle : des diagrammes
$E^2 \to E^1$ et $X^2 \to X_1$ s'y devinent, rien ne s'y lit. Nous ne
reproduisons pas ce décompte cellule par cellule : les carrés qui le portent
sont ceux transcrits ci-dessus, répétés avec des indices décalés.}

\page{19}
\begin{tikzcd}[column sep=large]
X_1 & X_2 \arrow[l, "p_0"', bend left=18] & X_3 \arrow[l, "p_{12}"', bend left=18]
\end{tikzcd}

avec de même $p_1 \colon X_2 \to X_1$ et $p_{23}, p_{31} \colon X_3 \to X_2$.

avec $\delta_0 \colon X_1 \to X_2$ et $\delta_1, \delta_2, \delta_3 \colon X_2 \to X_3$,
$\mathfrak{S}_2$, $\mathfrak{S}_3$, $\rho, \sigma_1, \sigma_2, \sigma_3$, et
$p_{31}$.

\struck{$\rho = \sigma$}

\[
  \mathfrak{S}_2 \ \Bigl|\ X_2 \twoheadrightarrow X_1 \times X_1, \qquad
  \underbrace{p_0\sigma}_{\|} = p_1, \quad (p_0,p_1) \text{ épi},
\]

avec $\delta_0$ remontant de $X_1$. \marginal{iso ????}

\[
  X_3 \xrightarrow{\ (p_{12},\,p_{13})\ \text{épi}\ } X_2 \times X_1
     \xrightarrow{\ (p_0,p_1)\times\mathrm{id}\ } X_1 \times X_1 \times X_1,
\]
\[
  X_3 \xrightarrow{\ \text{épi}\ (p_{12},\,p_{23})\ } X_2 \times_{X_1} X_2
     \xrightarrow{\ \text{épi}\ } X_2 \times_{X_1} (X_1 \times X_1),
\]

les deux facteurs du second terme étant marqués $(p_1)$ et $(p_0)$.

Deux carrés cartésiens :

\begin{tikzcd}[column sep=small, row sep=small]
X_3 \arrow[r, "p_1"] & X_2 \\
X_2 \arrow[u, "\delta_3"] & X_1 \arrow[l, "\delta_0"] \arrow[u, "\delta_0"']
\end{tikzcd}
\qquad
\begin{tikzcd}[column sep=small, row sep=small]
X_2 & X_3 \arrow[l, "p_{12}"'] \\
X_1 \arrow[u, "k_1"] & X_2 \arrow[l, "\delta_1"] \arrow[u, "k'_2"']
\end{tikzcd}

\note{L'auteur écrit « cart » au centre de chacun de ces deux carrés : c'est lui
qui les déclare cartésiens. Les mêmes carrés sont écrits une seconde fois sur la
page avec $E$ au lieu de $X$ et $\partial$ au lieu de $\delta$ — le passage d'un
système à l'autre est ce que la page cherche.}

\begin{tikzcd}[column sep=small, row sep=small]
E_2 \arrow[r, "\partial_2"] \arrow[d, "k_1"'] & E_3 \arrow[d, "k'_2"] \\
E_1 \arrow[r, "\partial_1"'] & E_2
\end{tikzcd}
\qquad
\begin{tikzcd}[column sep=small, row sep=small]
E_3 \arrow[r, "k_2"] \arrow[d, "k'_2"'] & E_2 \arrow[d, "k_1"] \\
E_2 \arrow[r, "k_1"'] & E_1
\end{tikzcd}

Supposons $X_1$ réduit à un élément, donc $X_2$ est un $\mathfrak{S}_2$-ens. à
pt fixe marqué $\delta_2$, $X_3$ un $\mathfrak{S}_3$-ens. à pt fixe marqué
$\delta_3$ ; \struck{\ill} $\delta_2$ et $k_2$ compatibles avec
\struck{les \ill} \uncertain{les pointations} par $\delta_2, \delta_3$, avec

\begin{tikzcd}[column sep=large]
X_2 \arrow[r, "\delta_2", bend left=18] & X_3 \arrow[l, "k_2", bend left=18]
\end{tikzcd}

\[
  k_2\delta_2 = \mathrm{id}_{X_2}.
\]

\note{Le bas de la page porte, à gauche, le système dual
$X_1 \rightleftarrows X_2 \rightleftarrows X_3$ étiqueté $\delta^*_1, k^*_1,
\delta^*_2, k^*_2$ avec $\mathfrak{S}_2$ et $\mathfrak{S}_3$ au-dessus, et
plusieurs figures griffonnées — segments à deux ou trois points marqués,
entourés — qui accompagnent visiblement le compte des points fixes. Elles ne
portent aucune lettre et nous ne les dessinons pas.}

\page{20}
\[
  \zeta = (x,y), \qquad
  d\sigma(\zeta) = \uncertain{(y,x,x)}, \qquad
  \rho^2 d(\zeta) = \uncertain{(y,x,y)}.
\]

\note{Trois lignes en haut d'une page par ailleurs vide, de la même encre que la
page 19 : le calcul en coordonnées y est repris une dernière fois. Les deux
triplets sont écrits petit et rapidement ; les lettres se ressemblent, d'où le
doute.}

\end{document}
